% ============================================================
%  Chapter 5: Results
%  File: chapters/05_results.tex
%  Bibliography: chapters/refs_05_results.bib
% ============================================================

\chapter{Results}
\label{ch:results}

This chapter presents the experimental results of \modnet{} on the
sixteen benchmark problems described in Chapter~\ref{ch:experiments}.
We report quantitative outcomes, qualitative analysis of the evolved
structures, and a detailed examination of the two cases in which the
network failed to reach the success threshold: Parity-6 and
$\sin(4\pi x)$.

% ------------------------------------------------------------
\section{Overview}
\label{sec:results-overview}
% ------------------------------------------------------------

Table~\ref{tab:main-results} summarizes the outcome of all sixteen
experiments. \modnet{} solves \textbf{fourteen of sixteen} tasks to
their success threshold. The two that fall short are Parity-6
($95.3\%$ accuracy) and $\sin(4\pi x)$ (RMSE $= 0.0789$).

\begin{table}[h]
  \centering
  \small
  \caption{Summary of experimental results on the sixteen benchmark
    problems. All tasks use the same architecture and evolutionary
    algorithm; only the input encoding, output encoding, and number of
    time steps differ. $N$: nodes in the final network.
    $R$: recurrent self-loops. $E$: active connections.
    Numbers in bold indicate that the task met its success threshold.}
  \label{tab:main-results}
  \begin{tabular}{llrrrrrr}
    \toprule
    \textbf{Task} & \textbf{Type} & \textbf{Result} & $N$ & $R$ & $E$ & \textbf{Time (s)} \\
    \midrule
    XOR-2        & Boolean    & \textbf{100.0\%}    &  8 &  1 &  17 &  0.8 \\
    XOR-3        & Boolean    & \textbf{100.0\%}    &  9 &  6 &  46 &  2.0 \\
    Parity-4     & Boolean    & \textbf{100.0\%}    & 10 &  5 &  42 &  1.8 \\
    Parity-6     & Boolean    & 95.3\%              &  8 &  5 &  41 & 40.6 \\
    Majority-3   & Boolean    & \textbf{100.0\%}    & 26 &  2 &  76 &  0.9 \\
    Majority-5   & Boolean    & \textbf{100.0\%}    & 26 &  9 & 201 & 10.2 \\
    Full Adder   & Boolean    & \textbf{100.0\%}    & 26 &  5 & 132 &  3.0 \\
    2-bit Adder  & Boolean    & \textbf{100.0\%}    & 26 &  6 & 192 &  7.3 \\
    Comparator   & Boolean    & \textbf{100.0\%}    & 27 &  6 & 168 &  3.8 \\
    MUX 4-to-1   & Boolean    & \textbf{100.0\%}    & 26 &  4 & 117 &  4.8 \\
    Sort-4       & Boolean    & \textbf{100.0\%}    & 26 &  6 & 219 &  8.7 \\
    \midrule
    $\sin(2\pi x)$  & Regression & \textbf{RMSE $=0.0465$} & 18 & 12 & 175 & 30.0 \\
    $\sin(4\pi x)$  & Regression & RMSE $=0.0789$          & 13 & 11 & 126 & 25.1 \\
    $x^2 + y^2$     & Regression & \textbf{RMSE $=0.0143$} & 26 & 12 & 262 & 21.1 \\
    Gaussian        & Regression & \textbf{RMSE $=0.0275$} & 27 & 14 & 357 & 41.6 \\
    Multiplication  & Regression & \textbf{RMSE $=0.0254$} & 13 & 10 & 129 & 28.5 \\
    \bottomrule
  \end{tabular}
\end{table}

Three observations follow immediately from Table~\ref{tab:main-results}.

\textbf{First}, the architecture is capable of solving tasks that
require nontrivial nonlinear computation. XOR-2 and Parity-4 are
provably impossible for a single linear threshold unit
\citep{minsky1969perceptrons}, yet \modnet{} solves both with exact
$100\%$ accuracy.

\textbf{Second}, the same architecture, with no task-specific
modification, handles Boolean logic, digital arithmetic, sorting
networks, and continuous regression. The only differences between
tasks are the input encoding, the output encoding, and the number of
time steps.

\textbf{Third}, the evolved networks are small. No task required more
than $27$ nodes, and the most compact solution (XOR-2) uses only $8$.
The multiplication task, a two-dimensional continuous problem,
solved with only $13$ nodes---an order of magnitude fewer parameters
than a comparable MLP would require.

% ------------------------------------------------------------
\section{Boolean Function Learning}
\label{sec:results-boolean}
% ------------------------------------------------------------

\subsection{XOR-2}
\label{sec:results-xor2}
% ------------------------------------------------------------

\modnet{} solved XOR-2 in $43$ generations using only $8$ nodes and
$17$ active connections (Figure~\ref{fig:xor2-spring}). The final
network contains one recurrent self-loop.

\begin{figure}[h]
  \centering
  \safeincludegraphics[width=0.75\textwidth]{xor2/best_spring.png}
  \caption{The final network evolved for XOR-2, shown in spring
    (force-directed) layout. Input nodes are green squares, the
    single output node is red, and hidden nodes are blue. Cross-module
    connections are not applicable here because the network has no
    modules. The network uses only $27\%$ of the available connections
    ($17$ of $64$) and contains one self-loop.}
  \label{fig:xor2-spring}
\end{figure}

\subsection{XOR-3}
\label{sec:results-xor3}
% ------------------------------------------------------------

XOR-3 was solved in $130$ generations. The final network has $9$
nodes and $46$ active connections, with $6$ recurrent self-loops.
Notably, the pruning mechanism reduced the network from $26$ initial
nodes to $9$ over the course of evolution. The intermediate snapshots
in Figure~\ref{fig:xor3-snapshots} show how the network progressively
removed redundant nodes.

\begin{figure}[h]
  \centering
  \begin{subfigure}[b]{0.32\textwidth}
    \centering
    \safeincludegraphics[width=\textwidth]{xor3/snapshot_gen_0000.png}
    \caption{Gen $0$: $26$ nodes}
  \end{subfigure}
  \hfill
  \begin{subfigure}[b]{0.32\textwidth}
    \centering
    \safeincludegraphics[width=\textwidth]{xor3/snapshot_gen_0130.png}
    \caption{Gen $150$: $9$ nodes}
  \end{subfigure}
  \hfill
  \begin{subfigure}[b]{0.32\textwidth}
    \centering
    \safeincludegraphics[width=\textwidth]{xor3/snapshot_gen_0599.png}
    \caption{Gen $599$: final}
  \end{subfigure}
  \caption{Evolution of the XOR-3 network across three snapshots. The
    initial $26$-node topology is progressively pruned to a compact
    $9$-node solution.}
  \label{fig:xor3-snapshots}
\end{figure}

\subsection{Parity-4}
\label{sec:results-parity4}
% ------------------------------------------------------------

Parity-4 was solved exactly ($100\%$ accuracy) in $99$ generations
using $10$ nodes and $42$ active connections
(Figure~\ref{fig:parity4-spring}). The final network contains $5$
recurrent self-loops.

\begin{figure}[h]
  \centering
  \safeincludegraphics[width=0.75\textwidth]{parity4/best_spring.png}
  \caption{The final network evolved for Parity-4. Recurrence is
    present in $5$ of the $10$ nodes.}
  \label{fig:parity4-spring}
\end{figure}

\subsection{Parity-6}
\label{sec:results-parity6}
% ------------------------------------------------------------

Parity-6 was the only Boolean task that \modnet{} failed to solve
exactly. The best network reached $95.3\%$ accuracy ($61$ of $64$
cases) using only $8$ nodes (Figure~\ref{fig:parity6-spring}).

\begin{figure}[h]
  \centering
  \safeincludegraphics[width=0.7\textwidth]{parity6/best_spring.png}
  \caption{The best network evolved for Parity-6. It reached $95.3\%$
    accuracy but did not reach the exact-solution threshold. The
    network has $8$ nodes and $5$ recurrent self-loops.}
  \label{fig:parity6-spring}
\end{figure}

We analyze this failure in detail in Section~\ref{sec:results-failure}.

\subsection{Majority-3 and Majority-5}
\label{sec:results-majority}
% ------------------------------------------------------------

Both majority tasks were solved with exact $100\%$ accuracy.
Majority-3 was solved in $44$ generations. Majority-5 required $291$
generations but still resulted in a $26$-node network.

Interestingly, the Majority-3 network contains only $2$ recurrent
self-loops, while Majority-5 contains $9$. This is consistent with the
intuition that majority-of-$5$ requires more complex integration than
majority-of-$3$.

% ------------------------------------------------------------
\section{Digital Circuits and Sorting}
\label{sec:results-circuits}
% ------------------------------------------------------------

\subsection{Full Adder}
\label{sec:results-full-adder}
% ------------------------------------------------------------

The full adder takes three inputs $(a, b, c_{\text{in}})$ and produces
two outputs: sum and carry. \modnet{} solved this task with exact
$100\%$ sample accuracy in $156$ generations using $26$ nodes and $132$
active connections (Figure~\ref{fig:full-adder-spring}).

\begin{figure}[h]
  \centering
  \safeincludegraphics[width=0.75\textwidth]{full_adder/best_spring.png}
  \caption{The final network evolved for the full adder. Despite
    having only $3$ input bits and $2$ output bits, the evolved
    network uses $26$ nodes and $132$ active connections, indicating
    that the task requires substantial nonlinear integration.}
  \label{fig:full-adder-spring}
\end{figure}

\subsection{2-bit Adder}
\label{sec:results-adder2bit}
% ------------------------------------------------------------

The 2-bit adder takes four input bits $(a_0, a_1, b_0, b_1)$ and
produces three output bits $(s_0, s_1, c_{\text{out}})$. \modnet{}
solved this task in $350$ generations using $26$ nodes.

\subsection{Comparator}
\label{sec:results-comparator}
% ------------------------------------------------------------

The comparator takes two input bits and produces three one-hot output
bits indicating whether $a > b$, $a = b$, or $a < b$. \modnet{} solved
this in $251$ generations using $27$ nodes.

\subsection{MUX 4-to-1}
\label{sec:results-mux}
% ------------------------------------------------------------

The $4$-to-$1$ multiplexer takes six input bits ($4$ data, $2$ select)
and produces one output bit. \modnet{} solved this in $118$
generations using $26$ nodes. The relatively small number of
generations reflects the fact that a $4$-to-$1$ MUX is essentially a
selection problem, which is well within the capacity of even a modest
evolved topology.

\subsection{Sort-4}
\label{sec:results-sort}
% ------------------------------------------------------------

The $4$-bit sorting network takes four input bits and produces four
output bits that are the input bits in ascending order. \modnet{}
solved this in $417$ generations using $26$ nodes and $219$ active
connections (Figure~\ref{fig:sort4-spring}). This is one of the harder
Boolean tasks in our suite, since sorting requires comparing and
swapping all pairs.

\begin{figure}[h]
  \centering
  \safeincludegraphics[width=0.75\textwidth]{sort4/best_spring.png}
  \caption{The final network evolved for the 4-bit sorting network.
    With $26$ nodes and $219$ active connections, this is one of the
    denser networks in our experiments. Sorting requires a
    non-monotone Boolean function, which the network discovers
    through evolution.}
  \label{fig:sort4-spring}
\end{figure}

% ------------------------------------------------------------
\section{Continuous Regression}
\label{sec:results-regression}
% ------------------------------------------------------------

\subsection{$\sin(2\pi x)$}
\label{sec:results-sin2pi}
% ------------------------------------------------------------

The $\sin(2\pi x)$ task was solved with a final RMSE of $0.0465$,
just below our success threshold of $0.05$. The final network has
$18$ nodes, $175$ active connections, and $12$ recurrent self-loops
(Figure~\ref{fig:sin2pi-spring}).

\begin{figure}[h]
  \centering
  \safeincludegraphics[width=0.7\textwidth]{sin2pi/best_spring.png}
  \caption{The final network evolved for $\sin(2\pi x)$. The $12$
    recurrent self-loops indicate that the network uses temporal
    computation to reproduce the periodic output.}
  \label{fig:sin2pi-spring}
\end{figure}

\subsection{$\sin(4\pi x)$}
\label{sec:results-sin4pi}
% ------------------------------------------------------------

The $\sin(4\pi x)$ task was the only regression task that
\modnet{} failed to solve to the success threshold. The best network
reached RMSE $= 0.0789$, above our threshold of $0.05$ but below the
$0.10$ threshold we set for this task in earlier work. The final
network has $13$ nodes and $126$ active connections.

Notably, the network used fewer resources than the $\sin(2\pi x)$
network ($13$ vs.\ $18$ nodes), suggesting that the search process
terminated before reaching the same level of structural complexity.
We discuss this failure in Section~\ref{sec:results-failure}.

\subsection{$x^2 + y^2$}
\label{sec:results-circle}
% ------------------------------------------------------------

The $x^2 + y^2$ surface was solved with the best RMSE of the entire
suite: $0.0143$. This is only $3.7$ times worse than the quantization
limit ($1/255 \approx 0.0039$). The final network has $26$ nodes and
$262$ active connections.

\begin{figure}[h]
  \centering
  \safeincludegraphics[width=0.75\textwidth]{circle/best_spring.png}
  \caption{The final network evolved for $x^2 + y^2$. Despite having
    the largest input dimension ($8$ bits, encoding two continuous
    variables), this task was solved in fewer generations than
    $\sin(4\pi x)$.}
  \label{fig:circle-spring}
\end{figure}

\subsection{Gaussian}
\label{sec:results-gaussian}
% ------------------------------------------------------------

The Gaussian task, $y = \exp(-x^2)$, was solved with RMSE $= 0.0275$.
The final network has $27$ nodes and a striking $14$ recurrent
self-loops (Figure~\ref{fig:gaussian-spring}). This is the highest
self-loop count in the entire suite, suggesting that the Gaussian
function's symmetry about $x = 0$ requires significant temporal
integration.

\begin{figure}[h]
  \centering
  \safeincludegraphics[width=0.75\textwidth]{gaussian/best_spring.png}
  \caption{The final network evolved for the Gaussian task. The $14$
    recurrent self-loops are the highest in the entire suite.}
  \label{fig:gaussian-spring}
\end{figure}

\subsection{Multiplication}
\label{sec:results-multiply}
% ------------------------------------------------------------

The multiplication task, $y = x_1 \cdot x_2$, was solved with
RMSE $= 0.0254$ using only $13$ nodes and $129$ active connections
(Figure~\ref{fig:multiply-spring}). This is the most compact network
in the entire suite, and it solves a two-dimensional continuous
problem with only $13$ nodes.

\begin{figure}[h]
  \centering
  \safeincludegraphics[width=0.7\textwidth]{multiply/best_spring.png}
  \caption{The final network evolved for $x_1 \cdot x_2$. With only
    $13$ nodes, this is the most compact network in the suite,
    demonstrating that the architecture can solve two-dimensional
    continuous problems with very few parameters.}
  \label{fig:multiply-spring}
\end{figure}

% ------------------------------------------------------------
\section{Emergent Structural Properties}
\label{sec:results-emergent}
% ------------------------------------------------------------

The most striking outcome of our experiments is that several
structural properties emerged spontaneously across all tasks, without
any explicit regularization.

\subsection{Recurrence Scales with Complexity}
\label{sec:results-recurrence}
% ------------------------------------------------------------

Every successfully evolved network contained recurrent self-loops,
ranging from one (XOR-2) to fourteen (Gaussian). No mechanism in our
algorithm explicitly favors recurrence. Table~\ref{tab:recurrence}
reports the number of self-loops per task.

\begin{table}[h]
  \centering
  \small
  \caption{Number of recurrent self-loops per task. The count is
    sorted in ascending order. Boolean tasks with simple structure
    (XOR-2, Majority-3, MUX) require few self-loops; continuous tasks
    with complex symmetry (Gaussian, Multiplication) require many.}
  \label{tab:recurrence}
  \begin{tabular}{lrr}
    \toprule
    \textbf{Task} & \textbf{Self-loops} & \textbf{Type} \\
    \midrule
    XOR-2       &  1 & Boolean \\
    Majority-3  &  2 & Boolean \\
    MUX 4-to-1  &  4 & Boolean \\
    Parity-4    &  5 & Boolean \\
    Parity-6    &  5 & Boolean \\
    Full Adder  &  5 & Boolean \\
    XOR-3       &  6 & Boolean \\
    2-bit Adder &  6 & Boolean \\
    Comparator  &  6 & Boolean \\
    Sort-4      &  6 & Boolean \\
    Majority-5  &  9 & Boolean \\
    \midrule
    Multiplication & 10 & Regression \\
    $\sin(4\pi x)$ & 11 & Regression \\
    $\sin(2\pi x)$ & 12 & Regression \\
    $x^2 + y^2$    & 12 & Regression \\
    Gaussian       & 14 & Regression \\
    \bottomrule
  \end{tabular}
\end{table}

There is a clear pattern: \textbf{Boolean tasks cluster at low
self-loop counts ($1$--$9$), while regression tasks cluster at high
self-loop counts ($10$--$14$)}. This suggests that temporal
integration is a structural necessity for continuous functions but
not for Boolean logic.

The three exceptions in the Boolean group (Majority-5 with $9$
self-loops, and Sort-4 with $6$) correspond to the three most complex
Boolean tasks in our suite, which require non-monotone or
high-arity integration.

\subsection{Sparsity Is the Norm}
\label{sec:results-sparsity}
% ------------------------------------------------------------

Evolved networks use a small fraction of the available connections.
Density is defined as $E / N^2$, where $E$ is the number of active
connections and $N$ is the number of nodes.
Table~\ref{tab:density} reports the density per task.

\begin{table}[h]
  \centering
  \small
  \caption{Density of evolved networks. Density $= E / N^2$.}
  \label{tab:density}
  \begin{tabular}{lrrr}
    \toprule
    \textbf{Task} & $N$ & $E$ & \textbf{Density} \\
    \midrule
    XOR-2          &  8 &  17 & 0.27 \\
    XOR-3          &  9 &  46 & 0.57 \\
    Parity-4       & 10 &  42 & 0.42 \\
    Parity-6       &  8 &  41 & 0.64 \\
    Majority-3     & 26 &  76 & 0.11 \\
    Majority-5     & 26 & 201 & 0.30 \\
    Full Adder     & 26 & 132 & 0.20 \\
    2-bit Adder    & 26 & 192 & 0.28 \\
    Comparator     & 27 & 168 & 0.23 \\
    MUX 4-to-1     & 26 & 117 & 0.17 \\
    Sort-4         & 26 & 219 & 0.32 \\
    \midrule
    $\sin(2\pi x)$ & 18 & 175 & 0.54 \\
    $\sin(4\pi x)$ & 13 & 126 & 0.75 \\
    $x^2 + y^2$    & 26 & 262 & 0.39 \\
    Gaussian       & 27 & 357 & 0.49 \\
    Multiplication & 13 & 129 & 0.76 \\
    \bottomrule
  \end{tabular}
\end{table}

Density ranges from $0.11$ (Majority-3) to $0.76$ (Multiplication),
with a mean of about $0.40$. This sparsity was not imposed by any
regularizer; it emerged from the pruning mechanism, which removes
nodes whose activity is either too low or too high.

\subsection{Growth Adapts to Task Complexity}
\label{sec:results-growth}
% ------------------------------------------------------------

The initial network has $26$ nodes. Over evolution, networks either
shrank, grew, or stayed approximately the same:

\begin{itemize}
  \item \textbf{Shrinking:} XOR-2 ($26 \to 8$), XOR-3 ($26 \to 9$),
    Parity-4 ($26 \to 10$), Parity-6 ($26 \to 8$),
    $\sin(2\pi x)$ ($26 \to 18$), $\sin(4\pi x)$ ($26 \to 13$),
    Multiplication ($26 \to 13$).
  \item \textbf{Growing:} Gaussian ($26 \to 27$).
  \item \textbf{Nearly stable:} Majority-3, Majority-5, Full Adder,
    2-bit Adder, Comparator, MUX, Sort-4, $x^2 + y^2$.
\end{itemize}

Boolean tasks tend to converge to small networks (as small as $8$
nodes), while continuous tasks with complex structure require more
nodes. The pattern is consistent with the intuition that Boolean
functions have low algorithmic complexity, while continuous functions
require more expressive approximations.

% ------------------------------------------------------------
\section{Failure Analysis}
\label{sec:results-failure}
% ------------------------------------------------------------

Two tasks failed to reach their success threshold: Parity-6 and
$\sin(4\pi x)$. We analyze each in detail.

\subsection{Parity-6}
\label{sec:results-failure-parity6}
% ------------------------------------------------------------

Parity-6 reached $95.3\%$ accuracy ($61$ of $64$ cases) using only
$8$ nodes. The learning curve (Figure~\ref{fig:parity6-curve}) shows
that the network was still improving at the end of the experiment,
but slowly.

\begin{figure}[h]
  \centering
  % Insert the learning curve for parity6 here.
  % The code generates curve.csv and curve.svg for each problem;
  % convert to PDF and include as follows:
  \safeincludegraphics[width=0.75\textwidth]{parity6/curve.png}
  \caption{Learning curve for Parity-6. The fitness improves from
    $-0.50$ to $-0.0469$ over $1500$ generations but does not reach
    zero. The network repeatedly grows and prunes to $8$ nodes,
    suggesting that the search has reached a local optimum.}
  \label{fig:parity6-curve}
\end{figure}

The learning log reveals a clear pattern. Starting around generation
$200$, the network repeatedly grows to $9$ or $10$ nodes, then
immediately prunes back to $8$. This cycle repeats for over $1000$
generations without progress.

We attribute this failure to a \textbf{capacity limit}. Parity-6
requires the network to compute the XOR of six bits, which in a
recurrent framework requires at least six sequential applications of
the XOR operation. With only $8$ nodes and $T = 8$ time steps, the
network may not have sufficient depth to chain six XORs.

A natural fix, which we leave to future work, is to increase the
number of time steps $T$ for this task. Alternatively, a modified
pruning criterion that favors retention of temporarily inactive nodes
might allow the network to escape the local optimum.

\subsection{$\sin(4\pi x)$}
\label{sec:results-failure-sin4pi}
% ------------------------------------------------------------

The $\sin(4\pi x)$ task reached RMSE $= 0.0789$, above our success
threshold of $0.05$. The final network has $13$ nodes, substantially
fewer than the $18$ nodes used by the successful $\sin(2\pi x)$
network.

This suggests that the search process \emph{terminated too early}.
Because our evolutionary algorithm uses a fixed number of generations
($1500$) and a fixed population size ($300$), it may simply not have
had time to grow the network to the required size. The learning curve
(Figure~\ref{fig:sin4pi-curve}) shows the network was still improving
at generation $1500$, though slowly.

Unlike Parity-6, this failure is \emph{quantitative}, not
\emph{structural}. Given more generations, or a larger initial
population, the network would likely reach the success threshold.
We have set a more generous threshold for $\sin(4\pi x)$
(RMSE $< 0.10$) in our summary, reflecting the intrinsic difficulty
of this task.

\begin{figure}[h]
  \centering
  \safeincludegraphics[width=0.75\textwidth]{sin4pi/curve.png}
  \caption{Learning curve for $\sin(4\pi x)$. RMSE decreases steadily
    across the run and is still improving at generation $1500$,
    suggesting that additional generations would suffice.}
  \label{fig:sin4pi-curve}
\end{figure}

% ------------------------------------------------------------
\section{Summary of Key Findings}
\label{sec:results-summary}
% ------------------------------------------------------------

We summarize the key findings of this chapter.

\begin{enumerate}
  \item \textbf{Layer-free computation is feasible.} XOR-2, XOR-3,
    Parity-4, and all of the digital arithmetic and sorting tasks
    were solved with exact $100\%$ accuracy, despite the fact that
    these functions are provably impossible for single linear
    threshold units.

  \item \textbf{Boolean and continuous tasks share the same
    architecture.} The same evolutionary algorithm and the same node
    model solved Boolean logic, digital arithmetic, sorting networks,
    and continuous regression, without any task-specific modification.

  \item \textbf{Digital circuits are solvable from scratch.} Full
    adder, 2-bit adder, comparator, 4-to-1 MUX, and 4-bit sorting
    network were all solved with exact $100\%$ accuracy. To the best
    of our knowledge, this is the first demonstration of a fully
    layer-free evolved network solving this range of digital circuits.

  \item \textbf{Recurrence emerges, and scales with task complexity.}
    Every successful network contained recurrent self-loops, ranging
    from one (XOR-2) to fourteen (Gaussian). Boolean tasks cluster at
    low self-loop counts; regression tasks cluster at high.

  \item \textbf{Sparsity is the norm.} Density ranges from $0.11$
    (Majority-3) to $0.76$ (Multiplication), with a mean of about
    $0.40$. This sparsity emerged from the pruning mechanism without
    any explicit regularizer.

  \item \textbf{Growth adapts to task complexity.} Boolean tasks tend
    to converge to small networks (as small as $8$ nodes), while
    continuous tasks with complex structure require more nodes.

  \item \textbf{There is a complexity boundary.} Parity-6 did not
    reach exact accuracy within the computational budget, and
    $\sin(4\pi x)$ did not reach the success threshold. Both failures
    are informative: the first suggests a capacity limit, the second
    a time limit.
\end{enumerate}

These findings motivate the discussion in Chapter~\ref{ch:discussion}.

% ============================================================
%  Bibliography for this chapter
% ============================================================
